\section{Experiments}
\label{sec: experiments}


\subsection{Datasets and Benchmarks}
\label{sec: datasets_benchmarks}
\begin{itemize}
    \item \textbf{NAVSIM v1} \cite{navsim-v1} filters OpenScene to remove near-trivial and erroneous scenes, reducing the train split size to 85k. During evaluation, NAVSIM v1 runs a non-reactive simulation at 10Hz for 4 seconds, then scores the driving agent with metrics including No at-fault Collision (NC), Drivable Area Compliance (DAC), Time To Collision (TTC), Comfort (Comf.), and Ego Progress (EP). These metrics are aggregated into the Predictive Driver Model Score (PDMS). We use the train split for training and the test split for evaluation.
    \item \textbf{NAVSIM v2} \cite{navsim-v2} improves simulation realism by enabling reactive traffic. It evaluates agents with the Extended Predictive Driver Model Score (EPDMS), adding metrics including Driving Direction Compliance (DDC), Traffic Light Compliance (TLC), Lane Keeping (LK), History Comfort (HC) and Extended Comfort (EC). 
\end{itemize}

\begin{table*}[t] \small
\centering
\caption{\textbf{Comparison with state-of-the-art methods on the NAVSIM v1.} 
NC: No at-fault Collision.
DAC: Drivable Area Compliance.
TTC: Time-To-Collision.
C.: Comfort.
EP: Ego Progress.
PDMS: Predictive Driver Model Score.
Abbreviations: 1x Cam (single front-view camera), Nx Cam (surround-view cameras), L (LiDAR).
 \dag: Using the best-of-N (N=6) strategy following~\cite{autovla}.
 }
 \vspace{-5mm}
\label{tab:sota_navsim_v1}
\setlength{\tabcolsep}{10pt}
\begin{center}
        \begin{tabular}{l|c|c|cc|ccc|>{\columncolor{gray!25}}c}
            \toprule
            \textbf{Method} & \textbf{Ref} & \textbf{Sensors} & \textbf{NC $\uparrow$} & \textbf{DAC $\uparrow$} & \textbf{TTC $\uparrow$} & \textbf{C. $\uparrow$} & \textbf{EP $\uparrow$} & \textbf{PDMS $\uparrow$} \\ 
            \midrule
            Human & - & - & 100 & 100 & 100 & 99.9 & 87.5 & 94.8 \\ 
            \midrule
            \multicolumn{9}{l}{\textit{BEV-based Methods}} \\
            UniAD~\citep{uniad} & CVPR'23 & 6x Cam & 97.8 & 91.9 & 92.9 & \textbf{100.0} & 78.8 & 83.4 \\
            TransFuser~\citep{transfuser} & TPAMI'23 & 3x Cam + L & 97.7 & 92.8 & 92.8 & \textbf{100.0} & 79.2 & 84.0 \\  
            PARA-Drive~\citep{para-drive} & CVPR'24 & 6x Cam& 97.9 & 92.4 & 93.0 & 99.8 & 79.3 & 84.0 \\  
            LAW~\citep{law} & ICLR'25 & 1x Cam & 96.4 & 95.4 & 88.7 & 99.9 & 81.7 & 84.6 \\
            Hydra-MDP~\citep{hydra-mdp} & arXiv'24 & 3x Cam + L & 98.3 & 96.0 & 94.6 & \textbf{100.0} & 78.7 & 86.5 \\
            DiffusionDrive~\citep{diffusiondrive} & CVPR'25 & 3x Cam + L & 98.2 & 96.2 & 94.7 & \textbf{100.0} & 82.2 & 88.1 \\
            WoTE~\citep{wote} & ICCV'25 & 3x Cam + L & 98.5 & 96.8 & 94.4 & 99.9 & 81.9 & 88.3 \\
            \midrule
            \multicolumn{9}{l}{\textit{VLA-based Methods}} \\
            AutoVLA~\citep{autovla} & NeurIPS'25 & 3x Cam & 98.4 & 95.6 & \textbf{98.0} & 99.9 & 81.9 & 89.1 \\
            ReCogDrive~\citep{recogdrive} & arXiv'25 & 3x Cam & 98.2 & \textbf{97.8} & 95.2 & 99.8 & 83.5 & 89.6 \\
            AutoVLA\dag~\citep{autovla} & NeurIPS'25 & 3x Cam & 99.1 & 97.1 & 97.1 & \textbf{100.0} & 87.6 & 92.1 \\
            DriveVLA-W0\dag & arXiv'25 & 1x Cam & \textbf{99.3} & 97.4 & 97.0 & 99.9 & \textbf{88.3} & \textbf{93.0} \\ 
            \midrule
            \rowcolor{gray!20}
            \model & - & 1x Cam & 98.9 & 95.3 & 96.1 & \textbf{100.0} & 81.6 & 87.9 \\ 
            \rowcolor{gray!20}
            \model\dag & - & 1x Cam & 99.2 & 96.6 & 96.9 & \textbf{100.0} & 83.4 & 89.8 \\ 
            \bottomrule
        \end{tabular}
\end{center}
\vspace{-9mm}
\end{table*}


\begin{table}[t] \small
\centering
\setlength{\tabcolsep}{16pt}
\caption{\textbf{Ablation of future image prediction.} 
PDMS: NAVSIM v1 \cite{navsim-v1} benchmark, EPDMS: NAVSIM v2 \cite{navsim-v2} benchmark. }
\vspace{-3mm}
\label{tab:ablation_frame}
        \begin{tabular}{l|c|c}
            \toprule
            \textbf{Setting} & \textbf{PDMS $\uparrow$} & \textbf{EPDMS  $\uparrow$}\\ 
            \midrule
            Random Frame & 77.3 & 75.2  \\ 
            Action Only & 51.8 & 48.9  \\ 
            Next Frame & \textbf{77.9} & \textbf{76.0} \\ 
            \bottomrule
        \end{tabular}
\end{table}


\begin{table}[!t]\small
\centering
\setlength{\tabcolsep}{16pt}
\vspace{-2mm}
\caption{\textbf{Ablation of action expert.} 
PDMS: NAVSIM v1 \cite{navsim-v1} benchmark, EPDMS: NAVSIM v2 \cite{navsim-v2} benchmark. }
\vspace{-3mm}
\label{tab:ablation_architecture}
        \begin{tabular}{l|c|c}
            \toprule
            \textbf{Setting} & \textbf{PDMS $\uparrow$} & \textbf{EPDMS  $\uparrow$}\\ 
            \midrule
            Use Diffusion & 19.7 & 19.6  \\ 
            Use VLM & 77.6 & 75.7  \\ 
            Action Expert & \textbf{77.9} & \textbf{76.0} \\ 
            \bottomrule
        \end{tabular}
\vspace{-7mm}
\end{table}



\subsection{Implementation Details}
\label{sec: implementation}
\textbf{Instruction Annotation}.
The driving instructions in our \textbf{InstructScene} dataset were generated by a fully-automated two-stage annotation pipeline. We select Qwen2.5-VL-72B-Instruct \cite{qwen2_5} as our annotation model for its powerful visual understanding capabilities. The inputs of each scene are 14 consecutive frames captured by the front-view camera at 2Hz, with a resolution of $(1920, 1080)$. The first 4 frames are considered past and current observations, and the last 10 frames are future observations that will not be available to the driving agent in the inference stage. 
\begin{itemize}
    \item \textbf{Stage One: Scene Understanding}. In stage one, we prompt the model with two requests, designed to convert the visual inputs and expected actions of the driving agent into natural language descriptions. We first instruct the model to describe the scene in the first 4 frames and to identify the traffic participants as well as the static objects. We then instruct it to describe the vehicle's driving behavior in the next 10 frames and its interaction with previously observed traffic participants. 
    \item \textbf{Stage Two: Instruction Formulation}. In stage two, we combine the visual inputs with the scene descriptions generated in stage one, and prompt the model to create concise driving instructions that would guide the driving agent to predict the actions described in stage one. 
\end{itemize}
Since VLMs struggle to accurately perceive ego-vehicle motion, we generate supplementary rule-based instructions. We classify scenes using speed, acceleration, and turn rate thresholds, converting them into natural language. While these closed-set instructions lack diversity, they provide precise ego-motion cues. We then use them as auxiliary prompts for the VLM, combining both to generate accurate and diverse driving instructions. With this pipeline, we annotated $85109$ scenes from the NAVSIM train split and $12144$ scenes from the NAVSIM test split. 

\textbf{Training.}
The model is trained on the NAVSIM train split for 200k steps using 8 H20 GPUs. We set the number of historical images to 4, and predict 8 future actions as well as the future image at the end of the actions. We set the learning rate to 2e-5 with 2500 warmup steps, and use a per-device batch size of 1. The weights of action and image loss are $\lambda_A = \lambda_V =1.0$. We also maintain an EMA model with a decay rate of 0.9999, which is saved as checkpoints. 


\subsection{Main Results}
\label{sec: main_results}
As shown in Tables~\ref{tab:sota_navsim_v2} and~\ref{tab:sota_navsim_v1}, our model demonstrates competitive performance on both NAVSIM benchmarks. On NAVSIM v2, it scores 86.9 EPDMS without any additional performance-enhancing techniques, which is comparable to SOTA. Using the best-of-N strategy as prior works~\cite{autovla, drivevla-w0}, it achieves top performance on NAVSIM v2, surpassing state-of-the-art methods on several metrics, including Driving Direction Compliance, Traffic Light Compliance, Lane Keeping, and History Comfort. These results suggest that \model has learned robust instruction following capabilities and benefited from future image prediction training. On NAVSIM v1, our model achieves 87.9 PDMS, matching multi-modal BEV methods, and improves to 89.8 with the best-of-N strategy. 
We note that \model achieves lower performance compared to state-of-the-art VLA-based methods on NAVSIM v1. This discrepancy is partially attributed to NAVSIM v1's inbalanced metrics, which disproportionately favor risk-averse policies over alternative, equally valid strategies learned by our model. 
Furthermore, competing VLA-based methods either require supplementary inputs such as multi-view images with high resolutions, or integrate CoT reasoning via additional RL training. 
Critically, these performance-enhancing mechanisms operate independently of our model’s core architecture and may be modularly incorporated without modifying the primary design. 

\begin{figure}[t]
\begin{center}
    \centering
    \includegraphics[width=\columnwidth]{figures/nVA.pdf}
    \vspace{-8mm}
    \captionof{figure}{\textbf{Ablation of interleaving image-action sequences}.
We compare the training losses of models trained on non-interleaving sequences (original) and interleaving sequences of different lengths.
}
\label{fig:nva}
\end{center}%
\vspace{-10.5mm}
\end{figure}

\begin{figure*}[t]
\centering
\includegraphics[width=1.0\textwidth]{figures/Additional_Examples.pdf}
\vspace{-6mm}
\caption{\textbf{Instruction-based planning examples}.
We visualize the effects of language instructions on action planning with front-view camera images and BEV maps. 
}
\label{fig:additional_examples}
\vspace{-7mm}
\end{figure*}

\subsection{Experimental Analysis}
\label{sec: ablation}
\textbf{Future Frame Prediction}. 
We ablate future frame prediction with three configurations. First, we always predict the next future frame. Second, we randomly sample one of the 8 future frames and specify the chosen index in the text prompt. Third, we remove the future frame prediction task altogether. All three configurations are trained on 8 H20 GPUs, with other hyperparameters identical to training in \Cref{sec: implementation}. The results, shown in \Cref{tab:ablation_frame}, indicate that the task of future frame prediction indeed improves the planning capabilities of the model, but the exact choice of future frame has limited impact on performance. 

\textbf{Interleaving Observation and Action}. 
In our original design, only past images are provided to the model as reference. We argue that interleaving image and action helps the model learn their dynamics, resulting in faster convergence and lower loss during training. Following~\cite{drivevla-w0}, we train the model with interleaving image-action sequences. We ablate the training image-action sequence length with 2, 4 and 6. During training, we interleave the original 4 past images with 3 past actions. Figure~\ref{fig:nva} reveals that although models trained on interleaving sequences suffer from higher loss in the initial stages of training, which can be attributed to the discrepancy between their training designs, they converge significantly faster than models without interleaving sequences, eventually surpassing the latter. In addition, although all interleaving models use the same sequence length during training, those with longer image-action sequences during training show lower losses. 

 
\textbf{Independent Action Module}. 
To validate our design of an additional action expert, we ablate it against using the existing VLM or diffusion modules for action planning. Although these alternatives reduce model size, they increase computational cost because of the high dimensions of these modules. As shown in \Cref{tab:ablation_architecture}, our action expert model slightly outperforms the VLM-module-based planner and significantly surpasses the diffusion-module-based one, confirming the effectiveness of our architecture.


\begin{figure*}[t]
\centering
\includegraphics[page=1,width=0.985\textwidth]{figures/Instruction_cropped.pdf}
\includegraphics[page=2,width=0.985\textwidth]{figures/Instruction_cropped.pdf}
\vspace{-2mm}
\caption{\textbf{Future image generation conditioned on instructions and actions.}
In the same scenario, given three sets of instructions, the model plans three action sequences and generates their respective future images. All action sequences follow their instructions and all images are consistent with their actions. 
}
\label{fig:instruction_action}
\vspace{-5mm}
\end{figure*}

\textbf{Visualizations}. 
In addition to Figure~\ref{fig:teaser}, we provide two more examples to visualize the effect of instructions on \model's ability to adjust the vehicle's speed according to user instructions in Figure~\ref{fig:additional_examples}. We test two instructions in each scene and plot the predicted trajectories for the next 8 frames on the front-view-camera image as well as the Bird's-Eye-View (BEV) map. In both scenes, our model successfully increased, decreased, or maintained speed to follow the instructions. 
We also offer a qualitative evaluation of our model's ability to align both its action planning and image generation with user instructions in Figure~\ref{fig:instruction_action}. We select critical scenes where there can be multiple possible courses of action, e.g. approaching the intersection, encountering another vehicle. For each scene, we give the model different sets of instructions, then generate future actions and images sequentially. We observe that \model is able to generate future actions and images that are consistent with the instructions, indicating that our world modeling framework has successfully helped the model learn the dynamics of the driving environment. 

\textbf{VLA Baseline}.
As a straightforward baseline for instructional driving, we extend Qwen-2.5-VL~\cite{qwen2_5} with a planning head to predict future actions based on language instructions. Despite being trained on the same dataset with instruction annotations as \model, this VLA model performs poorly, achieving only $\sim 60$ PDMS and often failing to generate instruction-consistent trajectories. We attribute this limitation to the sparse low-dimensional action supervision, which is insufficient to bridge high-dimensional visual-language inputs and low-level driving actions. This motivates us to explore dense visual supervision from future prediction to improve instruction-based planning.
